\begin{lemma}[Situation 7.3: stored theta in the disjoint case]
Fix a finite loopless multigraph with its individual edge identities, a
palette of $k$ colors, a distinguished edge $e_0$, an input coloring,
an exact terminal tree-sequence $T$, a generating coloring, a
positive-rung ETT, a hierarchy base, a hierarchy counterexample $C$, and
a \noderef{HierarchyMinimumCounterexampleContext} for these data.  Let
$P$ be a \noderef{HierarchyReplacementPresentation}, let
$\alpha,\beta$ be colors, and assume
\[
 2\le j:=\noderef{HierarchyPathNumber}(C.\mathit{hierarchy})
   <\ell:=\noderef{HierarchyFullTailLength}(C).
\]
Assume that every vertex has a missing color under every proper coloring
leaving exactly $e_0$ uncolored.  Let $d$ be a
\noderef{HierarchySituationThreeDisjointData} for
$P,\alpha,\beta$.  Put
\[
 \sigma=d.\mathit{switched.coloring},\qquad
 H=d.\mathit{switched.counterexample.hierarchy},\qquad
 q=H.\mathit{dividerCount},
\]
and
\[
 D=\noderef{HierarchyStoredColorsAt}(H,q).
\]
Assume only the additional hypothesis
\[
                         d.\theta\in D.                       \tag{S}
\]
In particular, the full two-alternative condition (9) is the one already
contained in $d$; neither of its alternatives is strengthened.  Then
these data are impossible.
\end{lemma}
\begin{proof}
Write $R=d.\mathit{switched}$,
$\mathcal E_\sigma=R.\mathit{corresponding.replacementETT}$, and
$g=R.\mathit{counterexample.good}$.  In addition to $j,\ell,H,q,D$ put
\[
 A=\noderef{HierarchyStarredStage}(H,q),\quad
 B=\noderef{HierarchyFullTailPrefix}(C,j-2),\quad
 X=\noderef{HierarchyFullTailPrefix}(C,j-1),
\]
\[
 Y=\noderef{HierarchyFullTailPrefix}(C,j),\qquad
 Z=\noderef{HierarchyFullTailPrefix}(C,\ell-1).
\]
Write $y_i$ for the last vertex of the exact prefix $T(y_i)$ and
$y_\ell$ for the last vertex of $T$.  The fields of
\noderef{HierarchySituationThreeDisjointData} give
\[
 \beta\in\noderef{TreeMissingColors}(\sigma,B),\qquad
 \beta\in\noderef{MissingColors}(G,\sigma,y_\ell),\qquad
 d.\theta\in\noderef{MissingColors}(G,\sigma,y_j),             \tag{1}
\]
together with \noderef{IsProperEdgeColoring},
\noderef{IsOneEdgeUncolored} at exactly $e_0$, and an exact
$\sigma$-\noderef{TashkinovAugmentationSequence} from $X$ to $T$.

Apply \noderef{HierarchyCounterexamplePenultimateElementary} to the six
fields of $R.\mathit{minimumContext}$, in their stated order: properness,
exactly-one-uncoloredness, the complete lower-rung package,
fewer-divider induction, rung-local lexicographic minimality, and the
$2n+11$ estimate.  This gives
\[
 \noderef{ColoredVertexSetElementary}
    (G,\sigma,\noderef{TreeSequenceVertexSet}(Z)).              \tag{2}
\]
We first derive the required order from the actual hierarchy $H$.  The
\(\mathit{coreInsideFirst}\) field of \noderef{TashkinovHierarchy} is an
entrywise exact inclusion of the base core in the first ordinary stage.
Each \(\mathit{nested}\) field is an entrywise exact inclusion of
successive ordinary stages, its matching \(\mathit{strict}\) field makes
that inclusion proper, and \(\mathit{terminalStage}\) identifies the last
stage with $T$.  If $q=0$, the zero branch of
\noderef{HierarchyStarredStage} identifies $A$ with the base core, so
\(\mathit{coreInsideFirst}\), followed by the ordinary-stage chain and
\(\mathit{terminalStage}\), places $A$ below $T$.  If $q>0$, the positive
branch identifies $A$ with the last ordinary preterminal stage, whose last
\(\mathit{nested}\) and \(\mathit{strict}\) clauses place it properly below
the terminal stage.  The \noderef{SameGoodHierarchy} field of $R$ is used
only to identify these exact replacement stages, including $A$ and $T$,
with the stages underlying the full-tail prefixes of $C$.  Expanding
\noderef{HierarchyFullTailLength} and
\noderef{HierarchyFullTailPrefix}, and using $2\le j<\ell$, now yields
\[
             A\preccurlyeq B\prec X\prec Y\preccurlyeq Z\prec T. \tag{3}
\]
Every relation in (3) is an entrywise
\noderef{TreeSequenceInitialSegment}; in particular it retains every
individual displayed edge identity.  The first relation is allowed to
be equality when $j=2$.

The vertex $y_j$ belongs to $Z$, while every vertex of $X$ is a distinct
displayed vertex of $Z$, by (3) and vertex injectivity in
\noderef{TreeSequence}.  From (1)--(2),
\[
 d.\theta\notin\noderef{TreeMissingColors}(\sigma,X).          \tag{4}
\]
At the last stored index, the definition of
\noderef{HierarchyStoredColorsAt} is the literal equality
\[
 D=D_n^\sigma\setminus\noderef{TreeMissingColors}(\sigma,A),
 \qquad
 D_n^\sigma=\noderef{TashkinovStoredColors}
    (\mathcal E_\sigma.\mathit{series.data},n),                \tag{5}
\]
where $n$ is the terminal rung.  Hence (S) gives both
$d.\theta\in D_n^\sigma$ and
$d.\theta\notin\noderef{TreeMissingColors}(\sigma,A)$.
Together with (1)--(4) and elementarity, this says that $y_j$ is the
first theta-missing vertex, so
\[
 \noderef{HierarchyFirstMissingPosition}(\sigma,T,d.\theta).
       \mathit{val}=|A|+j.                                    \tag{6}
\]

Because $d.\theta\in D$, the `twoColors' clause of
\noderef{GoodHierarchy} says that the reserved set
$g.\Gamma^q_{d.\theta}$ has cardinality two.  Choose and name its two
distinct members
\[
 \varepsilon _1,\varepsilon _2\in g.\Gamma^q_{d.\theta},
 \qquad \varepsilon _1\ne\varepsilon _2.                    \tag{7}
\]
The `avoidance' clause places both colors in
\noderef{HierarchyStageMissingColors}$(H,q)$ and outside
\noderef{HierarchySegmentColors}$(H,q,d.\theta)$.  At $q>0$ the
stage-missing set is the missing-color set of $A=T_{n,q}$; at $q=0$ it
is the missing-color set of the terminal ladder tree $T_n^\sigma$, an
exact prefix of $A$.  Thus
\[
 \varepsilon _1,\varepsilon _2
       \in\noderef{TreeMissingColors}(\sigma,A).               \tag{8}
\]
At the last stored index the next hierarchy stage is the terminal tree
$H.\mathit{stage}(q+1)=T$.  Equation (6) gives the first-missing cutoff
$|A|+j$, and $j<\ell$ makes this strictly smaller than $|T|$.  Thus the
minimum in the definition of the segment-color cutoff is exactly
$|A|+j$.  The displayed interval from $A$ through $Y$ is therefore
precisely the portion controlled by the avoidance clause.  Consequently
\[
 \sigma(T.\mathit{edge}(r))\notin
       \{\varepsilon _1,\varepsilon _2\}
       \quad\text{for }|A|\le r<|A|+j.                        \tag{9}
\]

Apply \noderef{HierarchyTailMissingColorExclusion} to
$\mathcal E_\sigma$, the base $R.\mathit{replacementBase}$, the
hierarchy $H$, elementarity (2), index $j$, and theta missing at $y_j$.
It says that theta is absent from the same exact interval except in the
single case
\[
 q=0,\quad A=T_n^\sigma,\quad r=|A|,\quad
 \Theta_n\in\{\mathrm{SE},\mathrm{RE}\},
\]
where the edge is the actual certified connecting edge $f_n$, its
connecting set is the singleton $\{f_n\}$, and
$\sigma(f_n)=d.\theta\in D_n^\sigma$.  This is a genuine possibility
because of (5); we retain it rather than discard it.  Equations (9) and
this precise exception are equation (12) of the paper.

We next extract the one consequence of condition (9) that is needed for
the first switch.  Define the exact used-tail set
\[
 U_j=\{c:\exists r\in\operatorname{Fin}|T|,
 |A|\le r<|A|+j,\ \sigma(T.\mathit{edge}(r))=c\}.
\]
If $d.\mathit{conditionNine}$ takes its outside alternative, it states
directly that $\beta\notin U_j$.  If it takes its stored alternative,
it states
\[
 \operatorname{Disjoint}(g.\Gamma^q_\beta\cup\{\beta\},U_j),
\]
and membership of beta in the singleton again gives
$\beta\notin U_j$.  Thus in both alternatives, and with no other use of
the split,
\[
 \sigma(T.\mathit{edge}(r))\ne\beta
       \quad(|A|\le r<|A|+j).                                \tag{10}
\]

The field $d.\mathit{betaOutsideGamma}$ says beta is outside the union
of all last reserved pairs.  Since epsilon1 belongs to the pair for
theta, $\varepsilon _1\ne\beta$.  Also (8) and (S) imply
$\varepsilon _1\ne d.\theta$: epsilon1 is missing on $A$, whereas a
stored theta is not.  Finally $\beta\ne d.\theta$: beta has a missing
witness in $B$, theta has its first missing witness at the distinct
vertex $y_j$, and equality of the colors would contradict elementarity
(2).  Put
\[
 i_\varepsilon=
   \noderef{HierarchyFirstMissingPosition}(\sigma,T,\varepsilon _1),
 \qquad
 i_\beta=\noderef{HierarchyFirstMissingPosition}(\sigma,T,\beta).
\]
Equations (1), (3), and (8) give
$i_\varepsilon.\mathit{val}\le |A|$ and
$i_\beta.\mathit{val}\le |B|$.  If beta is not missing on $A$, these
inequalities give the strict first-witness order
$i_\varepsilon<i_\beta$; if beta is missing on $A$, the conditional
order premise of \noderef{ChangeLemma} is not required.  In the former
case, its conditional unused-edge premise is the restriction of (9) to
the interval ending at $i_\beta$.  Both colors are missing on $Z$ by
(1), (3), and (8).  Finally, epsilon1 is missing on the last ordinary
stage when $q>0$, and on $T_n^\sigma$ when $q=0$, exactly the first or
second case of the Change Lemma.  The six minimum-context fields used
above and elementarity (2) therefore give
\[
 \noderef{TwoColorComponent}(G,\sigma,\varepsilon _1,\beta,
   \operatorname{inl}(T.\mathit{vertex}(i_\varepsilon)))
 =
 \noderef{TwoColorComponent}(G,\sigma,\varepsilon _1,\beta,
   \operatorname{inl}(T.\mathit{vertex}(i_\beta))).           \tag{11}
\]

Call the common component in (11) $Q_0$, and let $Q_\ell$ be the
epsilon1--beta component through $y_\ell$.  These components are
disjoint, including all coincident-witness and singleton cases.  If the
two first witnesses in (11) coincide, that vertex misses both affected
colors, so its vertex-node is isolated in
\noderef{TwoColorIncidenceGraph}; hence $Q_0$ is a singleton and cannot
contain the distinct terminal vertex $y_\ell$.  If the witnesses are
distinct, \noderef{MissingColorIncidenceEndpoint} gives incidence
degree at most one at those two vertices and at the beta-missing vertex
$y_\ell$.  They are three distinct vertices, the first two lying in
$Z$ and the last lying outside $Z$.  Were $Q_0=Q_\ell$,
\noderef{TwoColorComponentSpanningPath} would span that component by a
simple path, while \noderef{TwoColorIncidenceDegreeAtMostTwo} permits
only the two path ends to have degree at most one.  This is impossible.
If $Q_\ell$ itself is a singleton, it plainly cannot contain either
distinct root of $Q_0$.  In every case
\noderef{TwoColorComponentsDisjoint} gives
\[
                              Q_\ell\cap Q_0=\varnothing,      \tag{12}
\]
and the endpoint lemma makes $Q_\ell$, including its singleton case, a
\noderef{TwoColorPathComponent}.

Define the first terminal switch in the paper's orientation by
\[
 \mu _4=\noderef{TwoColorComponentSwap}
   (G,\sigma,\varepsilon _1,\beta,\operatorname{inl}(y_\ell)). \tag{13}
\]
It is proper by \noderef{TwoColorComponentSwapProper}; the defining
swap fixes $\mathsf{none}$, so exactly $e_0$ remains uncolored.
Disjointness (12) preserves epsilon1 and beta missing at their first
witnesses.  Since beta is missing at $y_\ell$ under sigma, epsilon1 is
missing there under $\mu _4$ after the swap; if $Q_\ell$ is a singleton,
both colors were already missing.  Therefore
\[
 \varepsilon _1\in\noderef{MissingColors}
    (G,\mu _4,T.\mathit{vertex}(i_\varepsilon))
      \cap\noderef{MissingColors}(G,\mu _4,y_\ell),
 \quad
 \beta\in\noderef{MissingColors}
    (G,\mu _4,T.\mathit{vertex}(i_\beta)).                    \tag{14}
\]

For later route transport we record the exact pointwise effect on $Z$.
If a vertex-node $v\in V(Z)$ lies outside $Q_\ell$, all incident
edge-nodes lie outside it and every incident color is fixed.  If it lies
inside $Q_\ell$, (12) makes it different from the first epsilon1 and
beta witnesses.  Elementarity (2) therefore makes both affected colors
present at $v$; properness makes their incident edge identities unique,
connectedness puts both edge-nodes in $Q_\ell$, and (13) exchanges their
values.  Expanding \noderef{PresentColors} and taking complements gives
\[
 \noderef{MissingColors}(G,\mu _4,v)
   =\noderef{MissingColors}(G,\sigma,v)\qquad(v\in V(Z)).      \tag{15}
\]
The singleton case uses the first alternative and is included.  By
(9)--(10), every exact displayed edge from $A$ through $Y$ is fixed by
the first switch.

The literal Cases-One-Two discriminator for the ordered pair
$(\varepsilon _1,\beta)$ follows from (8).  Apply
\noderef{StableChangeCasesOneTwo} to the properness and
one-uncoloredness just proved, the six minimum-context fields,
elementarity (2), the path $Q_\ell$, and separation (12).  Retain its
\noderef{HierarchyWeaklyStable} and \noderef{ModColoring} conclusions.

We now verify the extra data for the full reconstruction.  Apply
\noderef{TashkinovAugmentationSequenceSuffix} to the exact sigma-route
$d.\mathit{transportedSuffix}:X\longrightarrow T$ and the exact
relations $X\preccurlyeq Y\preccurlyeq T$ from (3).  Denote the
resulting exact suffix by
\[
                  \mathcal A^\sigma_{Y,T}:Y\longrightarrow T. \tag{16}
\]
Equations (14)--(15) put both affected colors in
\noderef{TreeMissingColors}$(\mu _4,Y)$.

For the source-to-$Y$ route, inspect the actual terminal positive
certificate of $\mathcal E_\sigma$ through
\noderef{TashkinovSeries} and \noderef{TashkinovSeriesStep}.  Let its
output be $C_\sigma$.  In PE, its \noderef{TashkinovClosure} gives a
sigma-augmentation route from $S_\sigma=T_n^\sigma$ to $C_\sigma$.
In SE or RE, it gives the individual edge $f_n$, the literal singleton
connecting-edge equation, the exact
\noderef{TreeSequenceAdjoinEdge} seed
$S_\sigma=T_n^\sigma+f_n$, and the route from that seed to $C_\sigma$.
The below-final field of \noderef{ExtendedTashkinovTree} and (3) give
$Y\preccurlyeq T\preccurlyeq C_\sigma$.

The source is also an exact prefix of $Y$.  Here $H$ is the actual
\noderef{TashkinovHierarchy} on $\mathcal E_\sigma$.  In PE at positive
divider count, its \(\mathit{baseStage}\) identity makes stage zero the
actual $T_n^\sigma$, and its iterated \(\mathit{nested}\) fields put that
tree below the last starred stage $A$; the same chain, ending with
\(\mathit{terminalStage}\), together with the exact stage identification
used in (3), gives
$T_n^\sigma\preccurlyeq A\preccurlyeq Y$.  In PE at divider zero, the
\(\mathit{parallel}\) field of \noderef{HierarchyBase} supplies the
\noderef{OrderedTreeUnion} defining $A$, whose
\(\mathit{leftSegment}\) field gives $T_n^\sigma\preccurlyeq A$; now
$H.\mathit{coreInsideFirst}$ followed by its ordinary
\(\mathit{nested}\) chain and \(\mathit{terminalStage}\), with the same
stage identification from (3), gives $A\preccurlyeq Y$.  In SE/RE at
positive divider count, $H.\mathit{baseStage}$ and the first
\(\mathit{strict}\) field say that the first ordinary stage has length at
least $|T_n^\sigma|+1$, the length of the exact seed.  The certified route
puts the seed below $C_\sigma$, while the \(\mathit{nested}\) chain and
\(\mathit{terminalStage}\), followed by $T\preccurlyeq C_\sigma$, put the
first ordinary stage below the same output.  Their entrywise comparison
therefore gives
$T_n^\sigma+f_n\preccurlyeq H.\mathit{stage}(1)
 \preccurlyeq A\preccurlyeq Y$.  In SE/RE at divider zero, the
nonparallel base clause gives $A=T_n^\sigma$, and the seed and $Y$ are
prefixes of the common output; $2\le j<\ell$ gives
$|A|+1\le |Y|$, so entrywise comparison gives the seed below $Y$.
Thus
\[
                  S_\sigma\preccurlyeq Y\preccurlyeq C_\sigma. \tag{17}
\]
Apply \noderef{TashkinovAugmentationSequenceRestrict} to the certified
route and the two relations (17).  Its conclusion is the actual
sigma-route $S_\sigma\longrightarrow Y$; this is the prefix half of the
route, distinct from the suffix (16).

Transport this restricted route to $\mu _4$.  At positive divider
count, \noderef{HierarchyWeaklyStable} contains a
\noderef{TreeColoringInvariant} on $A$.  At divider zero in PE, weak
stability is a \noderef{CombinedStableColoring}, and
\noderef{CombinedStableClosurePreserved} fixes and transports the
terminal ladder tree and base closure whose
\noderef{OrderedTreeUnion} is $A$.  At divider zero in SE/RE, the source
is the exact adjoined seed; (9)--(10) fix its extra edge, while the
\noderef{StableColoring} constituent of weak stability supplies the
vertexwise control on $T_n^\sigma$.  Combine that source control with
exact edge fixation through $Y$, and use (15) on $X$ together with
$X\prec Y$ and $|Y|=|X|+1$.  Every hypothesis of
\noderef{TashkinovAugmentationSequenceColoringTransport} now holds, and
it gives the exact $\mu _4$ route
\[
 \begin{cases}
 T_n^\sigma\longrightarrow Y,&\text{PE},\\
 T_n^\sigma+f_n\longrightarrow Y,&\text{SE or RE},
 \end{cases}                                                   \tag{18}
\]
retaining the same edge identity, exact seed, and singleton equation in
the second line.

For every $\eta\in D$, its first missing position is unchanged by the
first switch.  If $\eta=\beta$, this can occur only in the stored
alternative of condition (9), and (15) preserves its first witness in
$B\subseteq Z$.  The equality $\eta=\varepsilon _1$ is impossible by
(8) and the definition of $D$.  In particular theta is one of the
remaining stored colors and is globally unaffected by (13), using the
distinctness proved above; every other remaining stored color is also
unaffected.  Therefore
\[
 \noderef{HierarchyFirstMissingPosition}(\mu _4,T,\eta)
 =\noderef{HierarchyFirstMissingPosition}(\sigma,T,\eta)
 \qquad(\eta\in D).                                           \tag{19}
\]
For the Gamma cutoff, beta is outside the union of all last pairs.  If
epsilon1 belongs to $g.\Gamma^q_\eta$ for $\eta\in D$, then its
membership in $g.\Gamma^q_{d.\theta}$ and pairwise disjointness in
\noderef{GoodHierarchy} force $\eta=d.\theta$.  Equation (6) puts that
stored color's first missing position exactly at $y_j$, hence no later
than $Y$.  This proves the cutoff for both affected colors.

Relations (3), (12)--(19), the Cases-One-Two discriminator, properness,
and exact one-uncoloredness are all premises of
\noderef{StableChangeCasesOneTwoSplitTail} with split tree $Y$.  It
returns a certificate
\[
 \mathcal C_1:\noderef{CorrespondingExtendedTashkinovTree}
       (G,k,e_0,\mathcal E_\sigma,\mu _4,T),                   \tag{20}
\]
maximum property, a base $\mathcal B_1$, a hierarchy, a good witness
$g_1$, and
\[
                  s_1:\noderef{SameGoodHierarchy}(g,g_1).     \tag{21}
\]
By (14), the distinct vertices
$T.\mathit{vertex}(i_\varepsilon)\in V(A)$ and
$y_\ell\notin V(Z)$ both miss epsilon1 under $\mu _4$; hence $T$ is not
elementary.  The returned positive-rung, maximum, hierarchy, and good
data thus form a \noderef{HierarchyCounterexample} $C_1$ on
$\mathcal B_1$.  Apply
\noderef{HierarchyReplacementPresentationTrans} to
$R,\mu _4,\mathcal C_1,\mathcal B_1,C_1,s_1$.  It gives an actual
fixed-original presentation
\[
 R_1:\noderef{HierarchyReplacementPresentation}
   (G,k,e_0,\mathit{input},T,\phi,\mathit{ett},\mathit{base},C,
    \mathit{ctx}),\qquad R_1.\mathit{coloring}=\mu _4.        \tag{22}
\]
The cited transitivity theorem composes the intermediate correspondence
and transports the complete minimum context; no replacement-relative
certificate is silently treated as original-relative.

We now record equation (13) under $R_1$.  The same-good data (21)--(22)
identify the transported last stored set $D_1$, starred stage $A_1$,
and reserved pair at theta with $D,A$, and
$g.\Gamma^q_{d.\theta}$.  Hence theta remains stored and epsilon1 and
epsilon2 retain (7)--(9).  The first switch does not involve theta:
indeed theta differs from both affected colors by the distinctness
proved above.  Thus (4), (6), and the precise theta tail exclusion,
including its possible seed, remain true under $\mu _4$.  More
explicitly, (9)--(10) say that the first switch fixes every exact edge
from $A$ through $Y$, while the connecting-edge and extension-kind
agreement fields in the correspondence (20) identify the possible old
exceptional seed with the actual current SE/RE seed.  Equation (15)
supplies the same conclusion on all missing sets in $Z$.  Finally (14)
retains epsilon1 missing both at its hierarchy witness and at
$y_\ell$.  Consequently
\[
 \varepsilon _1,\varepsilon _2\notin
    \mu _4\langle Y-A\rangle,
 \qquad
 d.\theta\notin\noderef{TreeMissingColors}(\mu _4,X),         \tag{23}
\]
and theta is unused on the exact displayed interval $A$ through $Y$
except possibly on the actual divider-zero SE/RE connecting seed.  This
is the full formal equation (13), with its exceptional edge retained.

Apply \noderef{ChangeLemma} a second time, now to the ordered pair
$(\varepsilon _1,d.\theta)$ in the presentation $R_1$.  Epsilon1 has
its retained first witness in $A$, while (23) and theta missing at $y_j$
give
\[
 i_\varepsilon^{\mu _4}.\mathit{val}\le |A|
   <\noderef{HierarchyFirstMissingPosition}
       (\mu _4,T,d.\theta).\mathit{val}=|A|+j.                \tag{24}
\]
The conditional unused-edge premise is epsilon1's exact exclusion in
(23).  The colors are distinct, both are missing on $Z$, and epsilon1
retains its last-stage, or divider-zero terminal-ladder, missingness.
Thus the Change Lemma gives
\[
 \noderef{TwoColorComponent}(G,\mu _4,\varepsilon _1,d.\theta,
   \operatorname{inl}(T.\mathit{vertex}(i_\varepsilon^{\mu _4})))
 =
 \noderef{TwoColorComponent}(G,\mu _4,\varepsilon _1,d.\theta,
   \operatorname{inl}(y_j)).                                  \tag{25}
\]

Let $K_0$ be the common component in (25), and let $K_\ell$ be the
epsilon1--theta component through $y_\ell$.  The vertices at the first
epsilon1 witness, $y_j$, and $y_\ell$ are distinct by (3) and (24), and
they miss epsilon1, theta, and epsilon1 respectively.  The same
\noderef{MissingColorIncidenceEndpoint},
\noderef{TwoColorComponentSpanningPath}, and
\noderef{TwoColorIncidenceDegreeAtMostTwo} argument as above gives
\[
                              K_\ell\cap K_0=\varnothing.      \tag{26}
\]
The conclusion follows through
\noderef{TwoColorComponentsDisjoint}, includes a singleton
$K_\ell$, and the endpoint lemma makes $K_\ell$ a
\noderef{TwoColorPathComponent}.

Define the second terminal switch by
\[
 \mu _5=\noderef{TwoColorComponentSwap}
   (G,\mu _4,\varepsilon _1,d.\theta,
      \operatorname{inl}(y_\ell)).                            \tag{27}
\]
It is proper and leaves exactly $e_0$ uncolored as before.
Disjointness (26) preserves epsilon1 missing at its first witness and
theta missing at $y_j$.  Epsilon1 is missing at $y_\ell$ under
$\mu _4$, so theta becomes missing there under $\mu _5$; in the
singleton case both colors were already missing.  Thus
\[
 d.\theta\in\noderef{MissingColors}(G,\mu _5,y_j)
       \cap\noderef{MissingColors}(G,\mu _5,y_\ell).           \tag{28}
\]
The pointwise component argument used for (15), now with (26) and the
elementarity of $Z$ under $\mu _4$ supplied by (15), gives
\[
 \noderef{MissingColors}(G,\mu _5,v)
   =\noderef{MissingColors}(G,\mu _4,v)\qquad(v\in V(Z)).      \tag{29}
\]

Before reconstructing the ETT, apply
\noderef{StableChangeCasesOneTwo} to the current presentation $R_1$ and
the ordered pair $(\varepsilon _1,d.\theta)$.  The discriminator is
epsilon1's retained stage missingness.  Use $K_\ell$, separation (26),
properness, exact one-uncoloredness, the six fields of
$R_1.\mathit{minimumContext}$, and the penultimate elementarity obtained
from \noderef{HierarchyCounterexamplePenultimateElementary}.  Retain the
resulting weak-stability conclusion before invoking any reconstruction.

We use that conclusion to handle the edge deliberately left exceptional
in (23).  Let
$\mathcal E_1=R_1.\mathit{corresponding.replacementETT}$ and let
$T_n^1$ be its terminal ladder tree.  If $q=0$ and the terminal
certificate is SE or RE, its \noderef{TreeSequenceAdjoinEdge} datum
gives the exact seed $T_n^1+f_n^1$ and says that $f_n^1$ is incident
with $T_n^1$.  The nonparallel \noderef{HierarchyBase} clause identifies
$A_1=T_n^1$.  Same-good transport gives
$d.\theta\in D_1$, and the defining inclusion
$D_1\subseteq D_n^1$ follows from
\noderef{HierarchyStoredColorsAt}.  If
$\mu _4(f_n^1)=d.\theta$, the terminal-ladder
\noderef{StableColoring} constituent of
\noderef{HierarchyWeaklyStable} applies: the edge is
\noderef{TreeIncidentEdge} to $T_n^1$, and its old color theta lies in
$D_n^1$.  Its exact edge-value clause therefore proves
\[
                         \mu _5(f_n^1)=\mu _4(f_n^1).          \tag{30}
\]
If the seed does not have color theta, (23) says it does not have color
epsilon1, so the component swap fixes it directly.  Thus (30) holds in
all divider-zero SE/RE cases.  Every other displayed edge from $A_1$
through $Y$ is fixed directly by the two exclusions in (23).  This is
the paper's exceptional-seed argument: tail avoidance alone was not used
to fix $f_n^1$.

The old $\mu _4$ suffix for the second reconstruction is explicit.
Apply \noderef{TashkinovAugmentationSequenceComponentSwap} directly to
the named sigma-suffix $\mathcal A^\sigma_{Y,T}$ from (16) and the first
switch (13).  By (14)--(15), both epsilon1 and beta are tree-missing on
$Y$ under $\mu _4$, so the theorem transports the same exact chain and
all individual new-edge identities to a $\mu _4$-suffix
\[
                              Y\longrightarrow T.             \tag{31}
\]
Equations (28)--(29) put epsilon1 and theta in
\noderef{TreeMissingColors}$(\mu _5,Y)$.

For the source-to-$Y$ route, inspect the actual current terminal
certificate of $\mathcal E_1$ in
\noderef{TashkinovSeries}.  Let its output be $C_1^{\rm out}$ and its
source be $S_1=T_n^1$ in PE or the exact seed
$S_1=T_n^1+f_n^1$ in SE/RE, with the actual edge and singleton equation
from \noderef{TashkinovSeriesStep}.  Put
$H_1=R_1.\mathit{counterexample.hierarchy}$ and let $A_1$ be its last
starred stage.  The current \noderef{SameGoodHierarchy} certificate first
identifies the divider count and the exact tree $A_1$ with $q$ and $A$;
hence (3) gives
$A_1=A\preccurlyeq Y$.  We then use the fields of the actual
\noderef{TashkinovHierarchy} $H_1$, not those of the original hierarchy.
In PE at positive divider count, $H_1.\mathit{baseStage}$ identifies stage
zero with $T_n^1$, and its iterated \(\mathit{nested}\) fields put
$T_n^1\preccurlyeq A_1$.  In PE at divider zero, the current
\noderef{HierarchyBase} parallel clause supplies the
\noderef{OrderedTreeUnion} defining $A_1$, whose
\(\mathit{leftSegment}\) field gives $T_n^1\preccurlyeq A_1$; the
\(\mathit{coreInsideFirst}\), ordinary \(\mathit{nested}\), and
\(\mathit{terminalStage}\) fields of $H_1$ locate that same $A_1$ in the
current chain ending at $T$.  In SE/RE at positive divider count,
$H_1.\mathit{baseStage}$ and its first \(\mathit{strict}\) field make
$H_1.\mathit{stage}(1)$ at least $|T_n^1|+1=|S_1|$ entries long.  The
certified route puts the exact seed below $C_1^{\rm out}$, and the current
\(\mathit{nested}\) chain and \(\mathit{terminalStage}\), followed by the
below-final relation, put $H_1.\mathit{stage}(1)$ below that same output.
Entrywise comparison therefore puts $S_1$ below that first stage; iterating
the current \(\mathit{nested}\) fields through
$A_1=A\preccurlyeq Y$ gives $S_1\preccurlyeq Y$.  In SE/RE at divider
zero, the current nonparallel
base clause gives $A_1=T_n^1$; the seed and $Y$ are prefixes of the common
output, and the unchanged bound $2\le j<\ell$ gives
$|A_1|+1\le |Y|$, so the same entrywise length comparison gives
$S_1\preccurlyeq Y$.  Finally $Y\preccurlyeq T$ by (3), while
$H_1.\mathit{terminalStage}$ identifies the current last hierarchy stage
with $T$ and the below-final field puts $T\preccurlyeq C_1^{\rm out}$.
Thus
\[
                     S_1\preccurlyeq Y\preccurlyeq C_1^{\rm out}. \tag{32}
\]
Apply \noderef{TashkinovAugmentationSequenceRestrict} to this actual
$\mu _4$-route.  It gives the current route
$S_1\longrightarrow Y$; it does not produce, and is not used in place
of, the suffix (31).

Apply \noderef{TashkinovAugmentationSequenceColoringTransport} to that
restricted route.  At positive divider count, weak stability gives
invariance on $A_1$.  At divider zero in PE, combined stability, the
preserved closure, and its ordered union provide source control.  At
divider zero in SE/RE, terminal-ladder stable coloring provides the
vertex control and (30) provides exact control of the additional seed
edge.  All remaining exact edge values through $Y$ were fixed above,
and (29) on $X$, together with $X\prec Y$ and $|Y|=|X|+1$, gives the
predecessor missing-set equality.  The conclusion is the exact
$\mu _5$-route
\[
 \begin{cases}
 T_n^1\longrightarrow Y,&\text{PE},\\
 T_n^1+f_n^1\longrightarrow Y,&\text{SE or RE},
 \end{cases}                                                   \tag{33}
\]
with the current individual seed identity and singleton equation
retained in the second line.

For every $\eta\in D_1$, its first missing position is unchanged under
the second switch.  Epsilon1 is outside $D_1$ because it is missing on
$A_1$.  If $\eta=d.\theta$, equality (29) preserves its first witness
$y_j$; every other stored color is unaffected by (27).  Theta itself is
outside every last Gamma pair by the transported field
$d.\mathit{thetaOutsideGamma}$.  If epsilon1 belongs to
$g_1.\Gamma^{q_1}_\eta$, its retained membership in the pair indexed by
theta and pairwise disjointness force $\eta=d.\theta$, whose first
missing position is $y_j\in Y$.  These are exactly the first-position
and Gamma-cutoff premises of
\noderef{StableChangeCasesOneTwoSplitTail}.

All remaining premises are (23), (26)--(33), properness, exact
one-uncoloredness, the Cases-One-Two discriminator, and the old suffix
(31).  The split-tail theorem returns a correspondence certificate
$\mathcal C_2$ relative to $\mathcal E_1$, maximum property, a base
$\mathcal B_2$, a hierarchy, a good witness $g_2$, and
\[
 s_2:\noderef{SameGoodHierarchy}
        (R_1.\mathit{counterexample.good},g_2).                \tag{34}
\]
Equation (28) makes the terminal tree non-elementary under $\mu _5$.
Thus the returned positive-rung, maximum, hierarchy, and good data form
a \noderef{HierarchyCounterexample} $C_2$ on $\mathcal B_2$.  Apply
\noderef{HierarchyReplacementPresentationTrans} to
$R_1,\mu _5,\mathcal C_2,\mathcal B_2,C_2,s_2$.  It supplies a second
fixed-original presentation $R_2$ satisfying
\[
                         R_2.\mathit{coloring}=\mu _5.         \tag{35}
\]

The path bounds $2\le j<\ell$, presentation (35), and the two
theta-missing assertions (28) are precisely a witness that
\[
 j\in\noderef{HierarchyEligibleCommonMissingIndices}
   (G,k,e_0,\mathit{input},T,\phi,\mathit{ett},\mathit{base},C,
    \mathit{ctx}).                                             \tag{36}
\]
Apply \noderef{HierarchySituationThreeEligibleIndexBound} to the fixed
data, the two path bounds, the every-vertex-missing hypothesis, and
(36).  Its conclusion is
$j<\noderef{HierarchyPathNumber}(C.\mathit{hierarchy})=j$, a
contradiction.  This proves the claim.
\end{proof}
